\documentclass[journal]{IEEEtran}

\usepackage[T1]{fontenc}
\usepackage[utf8]{inputenc}
\usepackage[spanish,es-noquoting,es-nolists,es-noshorthands]{babel}
\usepackage{amsmath,amssymb}
\usepackage{booktabs}
\usepackage{cite}
\usepackage{url}
\usepackage[hidelinks]{hyperref}

\begin{document}

\title{Estructura frente a tamaño: selección de conjuntos críticos
en el grafo de exposición de DeFi}

\author{David Sánchez Lombana \quad Juan Sebastián Cuesta Montenegro\\[2pt]
{\small\texttt{davidh.sanchez@urosario.edu.co} \quad
\texttt{juans.cuesta@urosario.edu.co}}\\[2pt]
{\small Universidad del Rosario, Bogotá, Colombia}%
\thanks{Proyecto final del curso de Teoría de Grafos, 2026-II.
Profesor: Daniel Bojacá. Entrega inicial.}}

\pagestyle{plain}

\maketitle

\begin{abstract}
Los protocolos de finanzas descentralizadas (DeFi) se componen entre sí
reutilizando los mismos activos como colateral, lo que induce un grafo dirigido
y ponderado de exposición crediticia en el que la caída de un protocolo puede
propagarse a los demás. El estado del arte ya construyó ese grafo y lo modeló
de forma predictiva, pero selecciona los nodos a someter a estrés por su
\emph{tamaño}: el protocolo mayor, los cinco mayores, un sector completo.
Este proyecto plantea el problema complementario y estructural: dado un
presupuesto de $k$ protocolos, hallar el conjunto cuya remoción fragmenta al
máximo la red, y comparar los criterios estructurales de selección --conjunto
dominante conectado, puntos de articulación, intermediación e inmunización
espectral-- contra los criterios de tamaño empleados actualmente. El daño se
mide de dos formas, estructural y ponderada por exposición, porque diez vértices
concentran el 84{,}5\,\% de la exposición total y contar vértices invertiría el
orden de importancia.
El problema es NP-duro, por lo que el trabajo se centra en algoritmos de
aproximación y heurísticas evaluadas sobre datos reales: 314 instantáneas
semanales del grafo entre 2018 y 2024. Se presentan la motivación, la
formalización del problema y los objetivos.
\end{abstract}

\begin{IEEEkeywords}
Teoría de grafos, fragmentación de redes, conjunto dominante conectado,
puntos de articulación, robustez de redes, finanzas descentralizadas.
\end{IEEEkeywords}

\section{Introducción}

\IEEEPARstart{L}{a} propiedad que distingue a las finanzas descentralizadas
(DeFi) de la banca tradicional es la \emph{composabilidad}: cualquier protocolo
puede integrar a otro sin permiso previo, de modo que el activo que un protocolo
emite se convierte en el colateral que otro acepta \cite{capponi2023defi}. Un
depósito puede ser representado por un token, ese token depositado en un segundo
protocolo como garantía de un préstamo, y el resultado reutilizado en un tercero.
El efecto agregado es una red de dependencias en la que el valor de un protocolo
depende del valor de activos emitidos por otros, y en la que un choque sobre un
solo activo puede propagarse en cascada por toda la red \cite{shu2026dexposurefm}.

Esta estructura es, de forma natural, un grafo: los protocolos son los vértices y
las relaciones de exposición son las aristas dirigidas y ponderadas. La literatura
de redes financieras estableció hace más de una década que la topología de esa red
--y no solo el tamaño de sus participantes-- determina si un choque se amortigua o
se amplifica. El resultado central de Acemoglu \emph{et al.}
\cite{acemoglu2015systemic} es una transición de fase: una red densamente conectada
es más resiliente que una dispersa frente a choques pequeños, y más frágil frente a
choques grandes. Jackson y Pernoud \cite{jackson2020survey} organizan el campo, y
Bardoscia \emph{et al.} \cite{bardoscia2021physics} revisan el aparato de física de
redes que se le ha aplicado.

Hasta hace poco, aplicar ese aparato a DeFi era inviable porque el grafo no
existía: había que reconstruirlo desde los registros de las cadenas de bloques.
Eso cambió en noviembre de 2025, cuando Wu \emph{et al.} publicaron
\emph{DeXposure} \cite{wu2025dexposure}, el primer conjunto de datos a gran escala
de exposición crediticia entre protocolos, con 43{,}7 millones de registros que
cubren más de 4\,300 protocolos y 602 cadenas, junto con su código. Tres meses
después, el mismo grupo publicó un modelo fundacional de grafos temporales sobre
esos datos \cite{shu2026dexposurefm}.

Con el grafo disponible y su modelado predictivo ocupado, queda abierta una
pregunta que es estrictamente de teoría de grafos y que ninguno de los dos
trabajos aborda. Ambos caracterizan el grafo mediante métricas descriptivas
--grado, entropía, asortatividad, cercanía, agrupamiento-- y el segundo somete la
red a pruebas de estrés, pero en todos los casos \emph{los nodos que se golpean se
eligen por su tamaño}: el protocolo con mayor valor bloqueado, los cinco mayores,
o un sector completo. Ningún escenario selecciona los nodos por su posición
estructural en el grafo. La pregunta que este proyecto plantea es la inversa:

\begin{quote}
\emph{Dado un presupuesto de $k$ protocolos, ¿qué criterio de selección
identifica al conjunto cuya caída simultánea fragmenta más el sistema, y
coincide ese conjunto con los protocolos más grandes?}
\end{quote}

Formulada así, la pregunta pertenece a un problema clásico: el de la
\emph{fragmentación de redes} por remoción dirigida de vértices, cuya
contraparte constructiva es el \emph{conjunto dominante conectado}. Un conjunto dominante conectado (CDS) de un grafo es un subconjunto
de vértices que induce un subgrafo conexo y tal que todo vértice fuera de él es
adyacente a alguno de sus miembros; funciona como una columna vertebral de
comunicación. En redes inalámbricas de sensores, los protocolos de control de
topología construyen precisamente un CDS para responder a la pregunta ``qué nodos
puedo apagar sin desconectar la red''
\cite{wightman2010topology,wightman2008a3}. Invertida, esa misma construcción
\emph{sugiere} --en grafos ralos como este, aunque no como equivalencia
general-- qué nodos, al caer, fragmentan el sistema. La transferencia de ese
aparato a un grafo de naturaleza completamente distinta es el núcleo de este
proyecto.

Existe un precedente metodológico directo en otro dominio. Koç \emph{et al.}
\cite{koc2014topological} estudian las transiciones de fase en fallas en cascada
en redes eléctricas comparando tres métricas espectrales --radio espectral,
resistencia efectiva de grafo y conectividad algebraica-- contra el umbral crítico
de carga, y atacan la red mediante remoción dirigida por intermediación. Su
hallazgo es pertinente: la conectividad algebraica \cite{fiedler1973algebraic}
captura el comportamiento de la cascada con una correlación cercana al 90\,\%,
mientras que el radio espectral apenas alcanza el 50\,\%, porque las transiciones
en robustez frente a cascadas son de conectividad y no de proceso. Este proyecto
adopta ese diseño experimental y lo traslada al grafo de exposición de DeFi.

\section{Descripción del problema}

\subsection{Modelo}

Para cada instante $t$ se define el grafo dirigido y ponderado
\[
G_t = (V_t, E_t, w_t),
\]
donde $V_t$ es el conjunto de protocolos activos en $t$; una arista
$(u,v) \in E_t$ indica que el protocolo $u$ mantiene exposición crediticia frente
al protocolo $v$, es decir, que parte del valor que $u$ custodia depende de
activos emitidos o respaldados por $v$; y $w_t(u,v) \in \mathbb{R}_{>0}$ es la
magnitud de esa exposición. La relación no es simétrica: que $u$ dependa de $v$ no
implica lo recíproco, por lo que el grafo es genuinamente dirigido. Se denota por
$\tilde{G}_t$ su simetrización, empleada cuando la propiedad de interés es la
conectividad y no el sentido de la propagación.

\subsection{Caracterización empírica de la instancia}

El proyecto no trabaja sobre grafos sintéticos. La instancia es el panel semanal
publicado con \emph{DeXposure} \cite{wu2025dexposure}, cuyas dimensiones se
verificaron directamente sobre los datos y se resumen en la Tabla~\ref{tab:datos}.

\begin{table}[!t]
\renewcommand{\arraystretch}{1.25}
\caption{Dimensiones del panel de exposición inter-protocolo}
\label{tab:datos}
\centering
\begin{tabular}{ll}
\toprule
\textbf{Propiedad} & \textbf{Valor} \\
\midrule
Instantáneas semanales      & 314 \\
Cobertura temporal          & 2018-04-23 a 2024-04-29 \\
Registros de arista         & 1\,258\,297 \\
Vértices distintos          & 3\,917 \\
Pares dirigidos distintos   & 38\,527 \\
Vértices por instantánea (2022--2024) & 1\,083 -- 3\,043 \\
Reciprocidad binaria (2022--2024) & 0{,}16 -- 0{,}27 \\
Componentes conexas (2022--2024)  & 24 -- 99 \\
Componente gigante (2022--2024)   & 93{,}4\,\% -- 98{,}8\,\% de los vértices \\
\bottomrule
\end{tabular}
\end{table}

Las cuatro últimas filas se reportan sobre el período 2022--2024, en el que el
grafo ya alcanzó su tamaño maduro; antes de 2022 la red tiene menos de mil
vértices y sus valores no son comparables. Dos hechos de esa caracterización
condicionan el diseño. El primero es que el grafo \emph{no} es conexo: se
descompone en decenas de componentes, de las cuales una concentra entre el
93\,\% y el 99\,\% de los vértices y el resto son fragmentos diminutos. En consecuencia, la conectividad
algebraica del grafo completo es nula por definición y carece de valor
informativo; el análisis debe restringirse a la componente gigante, y el número de
componentes pasa a ser una variable de interés por derecho propio. El segundo es
que la reciprocidad es baja pero no despreciable, lo que obliga a justificar
explícitamente cada paso en que se simetrice el grafo.

\subsection{Formulación}

Sea $\tilde{G} = (V, E, w)$ la simetrización de una instantánea y sea
$\mathcal{C}(H)$ el tamaño de la componente conexa mayor de un grafo $H$. Para un
conjunto $S \subseteq V$, sea $\tilde{G} - S$ el grafo que resulta de eliminar los
vértices de $S$ y sus aristas incidentes.

\paragraph{Observación diagnóstica} Preguntar por el conjunto de cardinalidad
mínima que \emph{desconecta} la red no tiene interés en esta instancia. En la
componente gigante de la instantánea del 6 de marzo de 2023, 630 de sus 1\,870
vértices --el 33{,}7\,\%-- tienen grado uno y el grado mediano es 2; existen por
tanto puntos de articulación y la conectividad de vértices es
$\kappa(\tilde{G}) = 1$. El óptimo es un solo protocolo de la periferia y no
informa de nada. La red es rala y con una periferia masiva de hojas, y esa es una
caracterización útil del grafo, no un obstáculo: el problema pertinente es cuánto
daño causa un presupuesto dado, no cuán pequeño puede ser el corte.

\paragraph{Fragmentación con presupuesto} El problema que este trabajo aborda es,
dado un presupuesto de $k$ protocolos, hallar el conjunto que maximiza el daño:
\[
S^{\ast}_k = \arg\min_{S \subseteq V,\ |S| = k} \ D(\tilde{G} - S),
\]
donde $D$ es una medida de daño. Es el problema de \emph{fragmentación de red} o
\emph{ataque dirigido óptimo} \cite{albert2000error}, y es NP-duro
\cite{garey1979computers}; también lo es la construcción del conjunto dominante
conectado de cardinalidad mínima, para el cual Guha y Khuller
\cite{guha1998approximation} proporcionan algoritmos de aproximación con garantía
logarítmica. La intratabilidad es lo que lo hace interesante como proyecto de
teoría de grafos: obliga a implementar y comparar aproximaciones y heurísticas en
lugar de resolver el óptimo.

\subsection{Dos medidas de daño, y por qué no basta una}

La elección de $D$ no es un detalle de implementación. La medida estructural
inmediata es el tamaño de la componente gigante restante,
\[
D_{\text{est}}(\tilde{G} - S) = \mathcal{C}(\tilde{G} - S),
\]
que cuenta vértices. Pero el grafo es ponderado precisamente porque la exposición
tiene magnitud, y en esta red esa magnitud está extraordinariamente concentrada:
\textbf{los diez vértices de mayor fuerza acumulan el 84{,}5\,\% de la exposición
total}, y el 1\,\% superior el 88{,}5\,\%. Bajo $D_{\text{est}}$, un fragmento de
quinientos protocolos periféricos cuenta más que uno que contenga a los tres
mayores prestamistas, lo cual invierte el orden de importancia financiera.

Se define entonces la medida ponderada como la fracción de exposición que
sobrevive en la componente gigante:
\[
D_{\text{pond}}(\tilde{G} - S) \;=\;
\frac{\displaystyle\sum_{(u,v) \in E[\,\mathrm{CG}(\tilde{G} - S)\,]} w(u,v)}
     {\displaystyle\sum_{(u,v) \in E} w(u,v)},
\]
donde $\mathrm{CG}(\cdot)$ denota la componente gigante y $E[\cdot]$ las aristas
inducidas por ella. El trabajo reporta \textbf{ambas curvas de daño}, y toma
$D_{\text{pond}}$ como medida principal: es la que responde a la pregunta
financiera --cuánta exposición queda expuesta al contagio-- mientras que
$D_{\text{est}}$ conserva la lectura puramente topológica y permite comparar con
la literatura de robustez de redes \cite{albert2000error}.

\subsection{Criterios de selección en competencia}

El aporte del proyecto es comparar, sobre la misma instancia y el mismo
presupuesto $k$, criterios de selección procedentes de tradiciones distintas:

\begin{enumerate}
\item \textbf{Tamaño}, el criterio del estado del arte: los $k$ vértices de mayor
      fuerza (suma de pesos incidentes), que aproxima el valor bloqueado con el
      que se definen los escenarios de estrés en \cite{shu2026dexposurefm}.
\item \textbf{Centralidad de intermediación} \cite{freeman1977betweenness}, la
      estrategia de ataque empleada por Koç \emph{et al.}
      \cite{koc2014topological}.
\item \textbf{PageRank}, por ser el componente estructural del índice de
      importancia sistémica de \cite{shu2026dexposurefm}.
\item \textbf{Puntos de articulación}, el criterio de conectividad puro.
\item \textbf{Conjunto dominante conectado}, construido con el algoritmo de
      aproximación de Guha y Khuller \cite{guha1998approximation} y con la
      familia de protocolos de control de topología A3
      \cite{wightman2008a3,wightman2010topology}.
\item \textbf{Inmunización espectral}, remoción voraz que maximiza la caída del
      radio espectral \cite{chen2016node}.
\end{enumerate}

\subsection{Por qué los criterios de tamaño no bastan}

Una medición preliminar propia sobre las 174 instantáneas posteriores a enero de
2021 justifica la pregunta. Al regresar el logaritmo de cada invariante contra los
logaritmos del número de vértices y de aristas se obtiene que la resistencia
efectiva de grafo queda explicada por el tamaño con $R^2 = 1{,}000$ y el radio
espectral con $R^2 = 0{,}987$; la conectividad algebraica, en cambio, solo con
$R^2 = 0{,}455$. Es decir: dos de los tres invariantes espectrales clásicos son,
en esta red, una medida de tamaño con otro nombre, y solo el invariante de
conectividad conserva información propia. Esto coincide con el hallazgo de Koç
\emph{et al.} \cite{koc2014topological} por una vía independiente y refuerza la
decisión de organizar el proyecto alrededor de la conectividad y no del espectro
de la matriz de adyacencia.

\subsection{Alcance y limitaciones asumidas}

El proyecto se acota deliberadamente. No se propone construir el grafo, que ya
existe \cite{wu2025dexposure}, ni predecir su evolución, que ya se hace
\cite{shu2026dexposurefm}. El peso de arista de
$D_{\text{pond}}$ es el de exposición del propio panel, no el valor total
bloqueado, porque Luo \emph{et al.} \cite{luo2024piercing} demostraron que esa
métrica está inflada por doble conteo. Se asumen tres limitaciones desde el
planteamiento: la resolución temporal es semanal, mientras que las cascadas de
liquidación ocurren en minutos; solo hay tres episodios de cascada documentados
dentro de la ventana de datos, lo que impide cualquier diseño estadístico basado
en el número de eventos; y el grafo es endógeno al riesgo, pues los protocolos
eligen a quién exponerse y un protocolo puede ser central precisamente porque se
le considera seguro.

\subsection{Validación externa disponible}

Que tres episodios no permitan un diseño estadístico no significa que no
permitan nada: admiten un \emph{estudio de caso}, y son la única validación
externa a la que este trabajo puede aspirar. La pregunta es si el conjunto
$S^{\ast}_k$ calculado sobre la instantánea inmediatamente anterior a cada
episodio --Terra/Luna en mayo de 2022, FTX en noviembre de 2022 y la pérdida de
paridad de USDC en marzo de 2023-- contiene a los protocolos que efectivamente
resultaron afectados.

Dos precisiones metodológicas, ambas obligadas por lo que muestran los datos. La
primera es que la verdad de campo \textbf{no} puede ser la desaparición del
protocolo del grafo: al comparar las instantáneas previas y posteriores a cada
episodio se observa que solo 2 de 12 protocolos del ecosistema Terra dejan de
aparecer en las ocho semanas siguientes, 1 de 6 en el caso de FTX y 2 de 15 en
el de USDC. Los protocolos no se borran del registro; lo que colapsa es su
exposición. La verdad de campo será por tanto la caída de la fuerza del vértice
dentro de la ventana del evento, con un umbral declarado explícitamente. La
segunda es que la lista de protocolos afectados se curará a mano a partir de
informes públicos posteriores a cada episodio, y no por coincidencia de cadenas
de texto en el nombre: una búsqueda por patrón sobre los nombres devuelve falsos
positivos evidentes como \emph{Nexus Mutual}, \emph{NFTX} o \emph{Solana
Farm}. Esa curaduría es una decisión del equipo y quedará documentada junto con
los resultados.

\section{Objetivos}

\subsection{Objetivo general}

Determinar, mediante algoritmos de teoría de grafos y para un presupuesto dado
de protocolos, los conjuntos cuya remoción fragmenta en mayor medida el grafo de
exposición inter-protocolo de las finanzas descentralizadas, y evaluar si los
criterios estructurales de selección superan a los criterios basados en tamaño
que emplea actualmente el estado del arte.

\subsection{Objetivos específicos}

\begin{enumerate}
\item \textbf{Modelar} el panel semanal de exposición inter-protocolo como una
      familia de grafos dirigidos y ponderados, estableciendo de forma explícita
      el criterio de arista, la métrica de peso y el tratamiento de la
      componente gigante y de los fragmentos.

\item \textbf{Caracterizar} cada instantánea mediante invariantes de conectividad
      --número de componentes, tamaño de la componente gigante, conectividad
      algebraica, distribución de grados y concentración del peso de exposición--
      y documentar su evolución entre 2018 y 2024.

\item \textbf{Implementar} los seis criterios de selección de vértices descritos
      en la Sección~II-D, incluyendo el algoritmo de aproximación de Guha y
      Khuller para el conjunto dominante conectado y la remoción voraz por
      inmunización espectral.

\item \textbf{Comparar} experimentalmente los criterios mediante \emph{las dos}
      curvas de daño definidas en la Sección~II-D --la estructural
      $D_{\text{est}}$ y la ponderada por exposición $D_{\text{pond}}$-- en
      función del presupuesto de remoción $k$, sobre instantáneas representativas
      del panel.

\item \textbf{Cuantificar} el solapamiento entre los conjuntos que produce cada
      criterio, para establecer si el conjunto estructuralmente crítico coincide
      o no con el conjunto de los protocolos más grandes.

\item \textbf{Contrastar} los conjuntos críticos obtenidos sobre la instantánea
      previa a cada uno de los tres episodios documentados --Terra/Luna en mayo de
      2022, FTX en noviembre de 2022 y la pérdida de paridad de USDC en marzo de
      2023-- contra los protocolos efectivamente afectados, tomando como verdad de
      campo el colapso de la exposición del vértice y asumiendo el carácter
      anecdótico de una validación con tres casos.

\item \textbf{Analizar} el costo computacional de cada algoritmo en función del
      tamaño de la instancia y discutir su viabilidad para el monitoreo periódico
      de una red de miles de vértices.

\item \textbf{Desarrollar} una herramienta reproducible que reciba una
      instantánea del grafo y devuelva los conjuntos candidatos junto con las
      curvas de daño correspondientes.
\end{enumerate}

% \section*{Agradecimientos}

% Se agradece al profesor Pedro Mario Wightman Rojas, cuyo trabajo doctoral en
% control de topología \cite{wightman2010topology} proporciona el aparato
% algorítmico que este proyecto traslada a un grafo de naturaleza distinta.
\bibliographystyle{IEEEtran}
\bibliography{referencias}

\end{document}
